\section{Where the active ingredients may lie}\label{sec:ingredients}
\subsection{Component and dismantling evidence}
The foundational depression dismantling trial ($N=150$) detected no advantage of complete cognitive treatment over behavioural activation, with or without automatic-thought skills, at termination or six months \citep{jacobson1996component}.
A larger depression trial ($N=440$) established behavioural activation's non-inferiority to CBT on the Patient Health Questionnaire-9 (PHQ-9) at 12 months \citep{richards2016cobra}.
The modified intention-to-treat difference was $0.1$ PHQ-9 points (95\% CI $[-1.3,1.5]$), within the trial's prespecified 1.9-point margin; the per-protocol analysis supported the same inference.
Behavioural activation was delivered by junior mental health workers and CBT by psychological therapists, so this pragmatic trial compared treatment packages and workforces rather than isolating one component.
It supports a viable treatment without an explicit restructuring module, not identical effects or a particular learning mechanism.
In anxiety:
\begin{itemize}
\item In a 108-trial meta-analysis, cognitive therapy, exposure and their combinations were efficacious with ``no differential efficacy for any treatment components'' \citep{norton2007metaanalytic}.
\item In a full-factorial trial of internet CBT for social anxiety ($N=464$), adding psychoeducation or exposure improved post-treatment outcomes; adding cognitive restructuring or attention training did not show superiority \citep{sipka2025factorial}.
The psychoeducation and exposure advantages were no longer significant four months after treatment, as conditions including cognitive restructuring or attention training improved further.
Substantial attrition limits the follow-up interpretation.
The exposure component itself included recording and comparing expected and observed outcomes, so it did not exclude cognitive processing.
\item Adding cognitive restructuring to prolonged exposure did not enhance outcome in assault survivors with PTSD \citep{foa2005pecr}, nor in children with social anxiety \citep{demooij2023social}.
\item A component network meta-analysis in panic disorder (72 studies) singled out interoceptive exposure, not cognitive restructuring, as efficacy-enhancing \citep{pompoli2018dismantling}.
\item In GAD, the network meta-analysis authors concluded that pure cognitive restructuring was not a decisively active ingredient when considered alone \citep{papola2024gad}; the paper was later corrected for a duplicated cohort and an error in standard-error handling \citep{papola2024correction}.
\end{itemize}

\subsection{Evidence against the claim}
These findings are not uniform.
\begin{itemize}
\item Cognitive therapy beat exposure in social phobia in a meta-analysis of 20 direct comparisons ($N=1{,}308$), with no difference in PTSD, OCD or panic \citep{ougrin2011exposure}.
\item In a trauma trial, adding cognitive restructuring to exposure was associated with PTSD rates of 31\% versus 63--75\% without it at 6 months \citep{bryant2008cr}.
\item In-session cognitive restructuring correlates with outcome ($r=.35$), though from only four studies and correlationally \citep{ezawa2023restructuring}.
\item In a single-session trial with elevated public-speaking anxiety ($N=249$), restructuring before exposure and behavioural experiments both outperformed exposure without explicit expectancy processing \citep{johnco2025restructuring}.
The small, non-significant difference in anxiety reduction favoured behavioural experiments, which also produced greater expectancy change in a shorter session.
Follow-up was limited to one week.
\item The exposure--cognition split is blurred: exposure may itself work by changing the expectancy of harm \citep{hofmann2008extinction}.
\end{itemize}

\subsection{Why the component evidence is weak on both sides}
A review of 16 depression component studies found a small pooled advantage for full treatments over versions with one component removed ($g=0.21$, 95\% CI $[0.03,0.39]$) \citep{cuijpers2017components}.
The components differed, so this is not an estimate of cognitive restructuring's unique contribution.
Only one study had low risk of bias, and all but one lacked power to detect effects smaller than $g=0.55$.
Mediation evidence is correlational: cognitive change tracks symptom change however it is produced \citep{cristea2015dysfunctional,lorenzoluaces2015complicated}, cognitive mediation of symptom reduction was limited in a depression cohort \citep[median $r=.06$;][]{vittengl2014mediators}, and reviewers say no common or specific factor meets causal criteria \citep{kazdin2007mediators,cuijpers2019commonfactors}.
A null component result is thus weak evidence that restructuring adds nothing, and a positive correlation is weak evidence that it matters.

\subsection{The strongest defence: practice when calm}
A defender can answer that patients rehearse when calm rather than undertake effortful restructuring during peak distress.
We accept the force of this possibility, but successful rehearsal would not by itself identify practice or habit formation as the active ingredient.
It could improve retrieval, cognitive skill, expectancies or behavioural routines; these alternatives require measurement and component tests.
Repeated performance does not reliably produce outcome-insensitive habits in human laboratory tasks \citep{dewit2018failures}.
Exposure protocols emphasise learning through expectancy violation \citep{craske2014maximizing}.
The relevant questions are what is rehearsed, what changes, and whether the acquired response transfers to situations with high executive demands.
